% experimentacion
Implementamos en \textit{python}, con ayuda de las bibliotecas de \textit{scapy}, \textit{numpy} y el código provisto por la cátedra, una serie de funciones dedicadas al filtrado de paquetes ---por tipo de protocolo y tipo de dispersión---, detección de símbolos ---y su probabilidad de ocurrencia---, cálculo de la \textit{información} que cada uno provee, y ---en términos generales--- el cálculo de la \textit{entropía} de cada fuente. El Código \ref{funciones} muestra la implementación de estas últimas dos funciones\footnote{El código concerniente a la detección de paquetes y armado de la fuente de información nula es omitido, ya que es aquel provisto por la cátedra}. En el caso del cálculo de la \textit{entropía}, se aprovechó la sintáxis de comprensión de listas para definir los valores de la sumatoria.

\vspace{1em}
\lstinputlisting[language=pseudo, label=funciones, caption={Implementación del cálculo de \textit{información} y de \textit{entropía} en \textit{python}.}, captionpos=b]{files/src/.code/funciones.py}

\vspace{1em}
Para la experimentación, trabajaremos sobre el tráfico de tres redes diferentes:

\begin{itemize}
    \item \textit{red 1}: como observada por un host con acceso por conexión wi-fi, ubicado en un edificio céntrico de la ciudad.
    \item \textit{red 2}: como observada por un host con acceso por conexión ethernet, ubicado en una zona suburbana de la ciudad.
    \item \textit{red 3}: como observada por un host con conexión wi-fi ubicado en una casa en el norte de la provincia.
\end{itemize}

Podemos considerar la ubicación de cada uno de estos hosts como ubicaciones dentro de distintos anillos concéntricos, en términos de densidad poblacional, dentro del área metropolitana de Buenos Aires.

Para cada red, tomaremos muestras de $10.000$ tramas en un horario y día ---que suponemos--- es de mayor tráfico (a partir de las $18$ hs, un viernes), y controlaremos la interacción del host con la red. Esto es, las mediciones se realizarán bajo un contexto de \textit{reposo} ---minimizando la cantidad de paquetes envíados--- y bajo uso \textit{intensivo}, con el objetivo de medir la influencia del tráfico propio en la fuente resultante. 

También, mediremos el tiempo de recolección de los datos, en el contexto de \textit{reposo}, a fin de obtener una noción respecto al flujo de tramas por unidad de tiempo en cada una de las redes analizadas.



Además, realizaremos un análisis de las direcciones IP dentro de los paquetes ARP. Ya que este es un protocolo de capa de enlace, y local para cada una de las redes, esto nos permitirá observar las IPs propias de cada red con mayor detalle. A partir de las tramas capturadas, generaremos una segunda fuente de información de memoria nula, tal que cada símbolo sea del tipo $<\ SRC,\ DST >$. De este modo, de existir algun host \textit{distinguido}, es decir que sobresale en la cantidad de información que provee, lo podremos identificar a través de los símbolos a los que pertenece.
